% SOB 编码计算量优化效果 - 折线图
% 使用方法: % SOB 编码计算量优化效果 - 折线图
% 使用方法: % SOB 编码计算量优化效果 - 折线图
% 使用方法: % SOB 编码计算量优化效果 - 折线图
% 使用方法: \input{figures/sob_computation_chart.tex}
% 字体: 中文使用宋体（继承文档设置），英文使用 Times New Roman（newtxtext）

\usetikzlibrary{arrows.meta}
\begin{tikzpicture}
\begin{axis}[
    every node/.append style={font=\rmfamily},  % 确保使用衬线字体（Times/宋体）
    width=0.85\linewidth,
    height=5.5cm,
    xlabel={稀疏率 (\%)},
    ylabel={迭代次数},
    xmin=30, xmax=80,
    ymin=0, ymax=2200,
    xtick={37.5, 50, 62.5, 73.2},
    legend pos=north west,
    legend style={font=\small},
    grid=major,
    grid style={dashed, gray!30},
    axis lines=left,
    xlabel style={font=\small},
    ylabel style={font=\small},
    tick label style={font=\small},
]

% 朴素迭代 - 蓝色实线
\addplot[
    color=blue,
    mark=square*,
    thick,
    mark size=3pt,
] coordinates {
    (37.5, 400)
    (50.0, 1200)
    (62.5, 1600)
    (73.2, 2048)
};
\addlegendentry{朴素迭代}

% 优化迭代 - 红色虚线
\addplot[
    color=red,
    mark=triangle*,
    thick,
    dashed,
    mark size=3pt,
] coordinates {
    (37.5, 152)
    (50.0, 600)
    (62.5, 974)
    (73.2, 1433)
};
\addlegendentry{SOB 优化迭代}

\end{axis}
\end{tikzpicture}

% 字体: 中文使用宋体（继承文档设置），英文使用 Times New Roman（newtxtext）

\usetikzlibrary{arrows.meta}
\begin{tikzpicture}
\begin{axis}[
    every node/.append style={font=\rmfamily},  % 确保使用衬线字体（Times/宋体）
    width=0.85\linewidth,
    height=5.5cm,
    xlabel={稀疏率 (\%)},
    ylabel={迭代次数},
    xmin=30, xmax=80,
    ymin=0, ymax=2200,
    xtick={37.5, 50, 62.5, 73.2},
    legend pos=north west,
    legend style={font=\small},
    grid=major,
    grid style={dashed, gray!30},
    axis lines=left,
    xlabel style={font=\small},
    ylabel style={font=\small},
    tick label style={font=\small},
]

% 朴素迭代 - 蓝色实线
\addplot[
    color=blue,
    mark=square*,
    thick,
    mark size=3pt,
] coordinates {
    (37.5, 400)
    (50.0, 1200)
    (62.5, 1600)
    (73.2, 2048)
};
\addlegendentry{朴素迭代}

% 优化迭代 - 红色虚线
\addplot[
    color=red,
    mark=triangle*,
    thick,
    dashed,
    mark size=3pt,
] coordinates {
    (37.5, 152)
    (50.0, 600)
    (62.5, 974)
    (73.2, 1433)
};
\addlegendentry{SOB 优化迭代}

\end{axis}
\end{tikzpicture}

% 字体: 中文使用宋体（继承文档设置），英文使用 Times New Roman（newtxtext）

\usetikzlibrary{arrows.meta}
\begin{tikzpicture}
\begin{axis}[
    every node/.append style={font=\rmfamily},  % 确保使用衬线字体（Times/宋体）
    width=0.85\linewidth,
    height=5.5cm,
    xlabel={稀疏率 (\%)},
    ylabel={迭代次数},
    xmin=30, xmax=80,
    ymin=0, ymax=2200,
    xtick={37.5, 50, 62.5, 73.2},
    legend pos=north west,
    legend style={font=\small},
    grid=major,
    grid style={dashed, gray!30},
    axis lines=left,
    xlabel style={font=\small},
    ylabel style={font=\small},
    tick label style={font=\small},
]

% 朴素迭代 - 蓝色实线
\addplot[
    color=blue,
    mark=square*,
    thick,
    mark size=3pt,
] coordinates {
    (37.5, 400)
    (50.0, 1200)
    (62.5, 1600)
    (73.2, 2048)
};
\addlegendentry{朴素迭代}

% 优化迭代 - 红色虚线
\addplot[
    color=red,
    mark=triangle*,
    thick,
    dashed,
    mark size=3pt,
] coordinates {
    (37.5, 152)
    (50.0, 600)
    (62.5, 974)
    (73.2, 1433)
};
\addlegendentry{SOB 优化迭代}

\end{axis}
\end{tikzpicture}

% 字体: 中文使用宋体（继承文档设置），英文使用 Times New Roman（newtxtext）

\usetikzlibrary{arrows.meta}
\begin{tikzpicture}
\begin{axis}[
    every node/.append style={font=\rmfamily},  % 确保使用衬线字体（Times/宋体）
    width=0.85\linewidth,
    height=5.5cm,
    xlabel={稀疏率 (\%)},
    ylabel={迭代次数},
    xmin=30, xmax=80,
    ymin=0, ymax=2200,
    xtick={37.5, 50, 62.5, 73.2},
    legend pos=north west,
    legend style={font=\small},
    grid=major,
    grid style={dashed, gray!30},
    axis lines=left,
    xlabel style={font=\small},
    ylabel style={font=\small},
    tick label style={font=\small},
]

% 朴素迭代 - 蓝色实线
\addplot[
    color=blue,
    mark=square*,
    thick,
    mark size=3pt,
] coordinates {
    (37.5, 400)
    (50.0, 1200)
    (62.5, 1600)
    (73.2, 2048)
};
\addlegendentry{朴素迭代}

% 优化迭代 - 红色虚线
\addplot[
    color=red,
    mark=triangle*,
    thick,
    dashed,
    mark size=3pt,
] coordinates {
    (37.5, 152)
    (50.0, 600)
    (62.5, 974)
    (73.2, 1433)
};
\addlegendentry{SOB 优化迭代}

\end{axis}
\end{tikzpicture}
